\begin{helper}
Fix \(\eta,\varepsilon_1\) and \(n\).  Use the usual notation
\[
L=\log n,\quad
k_R=(1+\eta)\sqrt{nL},\quad
K=\noderef{TwoBiteNaturalIndependenceScale}(1+\eta,n),\quad k=(K:\mathbb R),
\]
\[
m_R=\frac n{L^2},\quad
M=\noderef{TwoBiteNaturalReducedVertexCount}(n),\quad m=(M:\mathbb R),
\qquad
p=\frac12\sqrt{\frac Ln},\qquad
t_1=\frac{\sqrt{nL}}{\log L}.
\]
Assume \(\varepsilon_1>0\), assume the base-threshold conclusions
\[
0<m,\quad 0\le p\le\frac12,\quad 1\le2pm,\quad
k_R\le k<k_R+1,\quad m\le m_R<m+1,
\]
together with \(K\le n\), \(k<2k_R\), \(\log L>0\), and
\[
\frac{2k}{t_1}+1\le5(1+\eta)\log L.
\]
If also
\[
n^{4\eta}\le\frac{\varepsilon_1}{2}k_R
\qquad\text{and}\qquad
2L^2\le\frac{\varepsilon_1}{8}k_R,
\]
then the base positivity and rounding conclusions hold, and moreover
\[
n^{4\eta}k\le\frac{\varepsilon_1}{2}k^2,\qquad
2kL^2\le\frac{\varepsilon_1}{8}k^2.
\]
\end{helper}
\begin{proof}
The base positivity and rounding conclusions are copied directly from the
base-threshold hypotheses.  Since \(K\) is a natural number, \(k=(K:\mathbb R)\)
is nonnegative.  Multiplying
\[
n^{4\eta}\le\frac{\varepsilon_1}{2}k_R
\]
by this nonnegative \(k\) gives
\[
n^{4\eta}k\le\frac{\varepsilon_1}{2}k_Rk.
\]
Because \(\varepsilon_1>0\) and \(k_R\le k\), the right side is at most
\((\varepsilon_1/2)k^2\).  This proves the first power comparison.  The second
is identical: multiply
\[
2L^2\le\frac{\varepsilon_1}{8}k_R
\]
by the nonnegative \(k\), then use \(k_R\le k\) and \(\varepsilon_1>0\) to get
\[
2kL^2\le\frac{\varepsilon_1}{8}k^2.
\]
\end{proof}
